Whether single-cell foundation models encode causal regulatory logic, rather than statistical
co-expression, remains unclear. We trained TopK sparse
autoencoders on residual-stream activations from every layer of Geneformer V2-316M and scGPT
whole-human, producing atlases of 82,525 and 24,527 features, released as interactive platforms. The features are biologically organized: 29--59\% annotate to standard ontologies and form
co-activation modules with coherent identity across the stack. Under a capacity-matched comparison they are not hidden from linear
decomposition---truncated SVD needs rank 250--403 to match their reconstruction quality---but give
a sparse parameterization supporting per-feature annotation and intervention. Ablating one feature disrupts the genes it fires on by 35--67$\times$,
though not the wider process it is annotated with. Against CRISPRi perturbation in four cell lines,
only 1 of 66 transcription factors shows target-specific responses. Benchmarking against differential expression on the same cells, established network inference and
random controls explains why: that assay measures 6,546 genes and the median perturbed factor has
one curated target among them, so no method can demonstrate specificity. On perturbation data covering
the transcriptome the features do recover targets, for 15\% of factors against 1\% for
perturbations of non-regulatory genes. Sparse features therefore carry a small amount of genuinely
factor-specific regulatory information, detectable only where the assay measures the targets.
